\documentclass[conference]{IEEEtran}
\usepackage{cite}
\usepackage{amsmath,amssymb,amsfonts}
\usepackage{graphicx}
\usepackage{textcomp}
\usepackage{xcolor}
\usepackage{booktabs}

\begin{document}

\title{Multi-Recognizer Multi-Attacker Adversarial Training for Robust Face Recognition}

\author{
\IEEEauthorblockN{Jerome Jodie Harianto}
\IEEEauthorblockA{
ERG4902 Capstone Paper Submission\\
CUHK-Shenzhen 123010005@link.cuhk.edu.cn\\
Repository: https://github.com/ZeroOreos/FYP
}
}

\maketitle

\begin{abstract}
Face recognition systems are used for identity verification and identification, 
and are assessed through a suite of protocols ranging from clean benchmarks to 
robustness-oriented evaluations. Practical use of face recognition models 
requires real-world robustness considerations, where appearance variation, dirty-data 
conditions, and adversarial inputs require identity embeddings to remain both 
separable and stable. Recent adversarial ensemble training efforts have attempted 
to address part of this problem by increasing attack-side pressure during training. 
Following such efforts, this work investigates whether adding recognizer diversity 
improves robustness beyond attack diversity alone. To test this, a single ArcFace/iResNet100 target is trained on WebFace4M 
under five comparison conditions: clean training, recognizer-only pressure, single-attack 
training, attack-ensemble training, and full recognizer-plus-attacker training. The 
central comparison holds the attacker suite fixed to PGD, BPFA, and DFANet, then adds 
auxiliary ArcFace, CosFace, and CurricularFace recognizer heads during training only. 
This raises average robust accuracy from 52.32\% to 56.55\% and reduces average attack 
success from 44.32\% to 39.45\%, while verification accuracy decreases from 98.33\% to 
98.21\% and EER increases from 1.78\% to 1.93\%. Inference cost is unchanged because only 
the target recognizer is deployed, but training cost increases substantially. The result 
is a calibrated tradeoff: recognizer diversity adds measured digital robustness beyond 
attack diversity alone, but moves the model along a robustness–verification–training-cost 
frontier.
\end{abstract}

\begin{IEEEkeywords}
face recognition, adversarial training, ensemble robustness, ArcFace,
verification, biometric security
\end{IEEEkeywords}

\section{Introduction}
Face recognition (FR) systems are used for identity verification and
identification, and are commonly assessed through clean benchmarks such as LFW,
MegaFace, and IJB-C \cite{lfw,megaface,ijbc}. Deep models and margin-based
objectives, including DeepFace, FaceNet, VGGFace, SphereFace, and ArcFace, have
made these protocols strong references for clean recognition performance
\cite{deepface,facenet,vggface,sphereface,arcface}. Clean accuracy, however,
does not fully describe deployment reliability. FR systems may encounter
distribution shift, appearance variation, common corruptions, and deliberate
adversarial inputs. This has led to dirty-data and robustness-oriented
benchmarks, including corrupted or appearance-varied face evaluations and
fine-grained adversarial test frameworks \cite{facesec,oodface}.
Decision-based black-box attacks have also been demonstrated against FR systems
\cite{dongattack}, and surveys document diverse attack and defense settings
\cite{frsurvey}. For biometric systems, these failures matter because the same
embedding must support both high identity separability and stable decisions
under non-clean inputs.

Adversarial training is one way to improve reliability under such pressure,
but robustness studies repeatedly show that results depend on the training and
evaluation threat model. Robust optimization frames defense as a min--max
problem, ensemble adversarial training reduces dependence on a single
adversary, and evaluation suites such as AutoAttack, RobustBench, and
MultiRobustBench emphasize that narrow attack coverage can overestimate
robustness \cite{madry,tramer,autoattack,robustbench,multirobustbench}. In FR,
this concern is tied to verification: the model is not only predicting a
training identity, but learning an embedding space whose pairwise distances are
thresholded at test time. A training signal that improves adversarial
classification can still alter the clean verification geometry.

This suggests two sources of training pressure. Attack diversity determines
which perturbation procedures the target is exposed to. Recognizer-objective
diversity determines which embedding geometries judge the target during
training. Modern margin and quality-aware objectives show that changing the
recognition loss changes the difficulty structure of the embedding space
\cite{arcface,cosface,curricularface,magface,adaface,elasticface}. The
question is therefore not only whether a stronger attacker helps, but whether
recognizer-side diversity adds robustness once attack diversity is already
present. The same added pressure may also reduce verification quality or
increase training cost.

The practical setting matters because large-scale FR is already resource
intensive. RetinaFace-style alignment, WebFace-scale data, Partial FC, and
sub-center ArcFace address alignment, many identities, classifier-head memory,
and noisy web labels \cite{retinaface,webface260m,partialfc,subcenter}. Adding
online attacks and auxiliary recognizer objectives compounds these costs. For
that reason, this paper treats robustness as a tradeoff frontier rather than a
free improvement: robust accuracy, verification accuracy, equal error rate,
memory pressure, batch size, and wall-clock cost must be read together.

This work studies whether recognizer-objective diversity adds measurable value
after attacker diversity is already present. The experiments use one deployed
target recognizer and compare clean training, recognizer-only pressure,
single-attack training, attack-ensemble training, and the full
recognizer-plus-attacker setting. The goal is not to claim state-of-the-art FR,
but to isolate a specific design effect: whether training-time recognizer
pressure can improve digital adversarial robustness beyond a multi-attack
baseline, and what that improvement costs in verification quality and training
efficiency.

\section{Methods}
\subsection{Design Rationale}
The design separates the deployed recognizer from the training pressures used
to shape it. The deployed model is a single ArcFace/iResNet100 target. During
training, the target sees adversarial samples from several attack families and
auxiliary recognition objectives that evaluate the same clean or perturbed
faces. This keeps inference cost unchanged while testing whether the target can
absorb useful pressure from a broader training environment.

The sets are intentionally compact. ArcFace supplies the deployed angular
margin target, CosFace supplies a cosine-margin alternative, and CurricularFace
changes the weighting of hard examples. PGD supplies direct optimization,
BPFA supplies transfer-oriented pressure, and DFANet supplies FR feature-level
pressure. The full condition is therefore one recognizer trained under rotating
attacker pressure and training-only recognizer pressure, not a deployed
multi-model ensemble.

\subsection{Training Objective}
The defended target is an ArcFace model with an iResNet100 backbone and a
512-dimensional embedding. Training uses WebFace4M manifests with
RetinaFace-style aligned $112{\times}112$ crops. The configured subset contains
205,849 identities, selected with a fixed random seed and a minimum of two
images per identity, so all comparison conditions share the same identity
universe.

Let $f_{\theta}(x)\in\mathbb{R}^{512}$ denote the target embedding and
$g_{\theta}(x)=f_{\theta}(x)/\|f_{\theta}(x)\|_2$ its normalized form. A
margin head with normalized class weights $W$ produces logits of the form
\begin{equation}
    \ell_y(x)=s\,\phi_y(W_y^\top g_{\theta}(x)),\quad
    \ell_j(x)=s\,W_j^\top g_{\theta}(x),\ j\neq y ,
\end{equation}
where $\phi_y(\cdot)$ is the margin rule applied to the ground-truth class.
The auxiliary recognizers differ mainly through $\phi_y$, giving the same
samples different margin geometries rather than redundant copies of the target
classifier.

Adversarial examples are generated by a family-conditioned inner problem
\begin{equation}
    x^{adv}_{k} \approx
    \arg\max_{\|\delta\|_{\infty}\leq\epsilon_k}
    \mathcal{J}_{k}(x+\delta,y;\theta,\mathcal{S}),
\end{equation}
where $k\in\{\mathrm{PGD},\mathrm{BPFA},\mathrm{DFANet}\}$ and $\mathcal{S}$
is the available surrogate set. The optimized loss combines the pressure
terms separated by the ablation:
\begin{equation}
\mathcal{L} =
\lambda_c\mathcal{L}_{c} +
\lambda_a\mathcal{L}_{a} +
\lambda_r\mathcal{L}_{r} +
\lambda_s\mathcal{L}_{sim} +
\lambda_h\mathcal{L}_{hard}.
\end{equation}
Here $\mathcal{L}_{c}=\mathrm{CE}(\ell(x),y)$ and
$\mathcal{L}_{a}=\mathrm{CE}(\ell(x^{adv}),y)$ keep the target discriminative,
while
\begin{equation}
    \mathcal{L}_{r}=\sum_{r\in\mathcal{R}}\alpha_r
    \left[\mathrm{CE}(\ell_r(x),y)+\mathrm{CE}(\ell_r(x^{adv}),y)\right]
\end{equation}
adds pressure from recognizer heads
$\mathcal{R}=\{\mathrm{ArcFace},\mathrm{CosFace},\mathrm{CurricularFace}\}$.
The consistency term discourages excessive clean/adversarial embedding drift,
\begin{equation}
    \mathcal{L}_{sim}=1-\cos(g_{\theta}(x),g_{\theta}(x^{adv})),
\end{equation}
and $\mathcal{L}_{hard}$ emphasizes mined hard pairs or hard samples. The
recognizer ensemble affects optimization only. After training, only
$f_{\theta}$ is deployed.

\subsection{Training Pipeline}
Training follows a three-stage curriculum over 24 epochs. Clean warmup lasts
3 epochs, shallow adversarial pressure lasts 9 epochs, and full adversarial
pressure lasts 12 epochs. Shallow and full stages use the same attack families
with different strengths. For stage $t$, each batch is split as
\begin{equation}
    B_t=B^{clean}_t\cup B^{adv}_t,\quad
    |B^{adv}_t|=\rho_t |B_t|,
\end{equation}
where $\rho_t=0$ during clean warmup, increases during shallow training, and is
0.50 in the full condition. One primary attacker family is selected in
round-robin order for each adversarial batch. Impersonation and dodging intents
are balanced to cover the main FR failure modes.

Each iteration has a target branch and pressure branches. Clean images update
the ArcFace/iResNet100 target through the margin objective. During adversarial
stages, part of the batch enters the attack branch, where the selected attacker
is optimized against the current target or surrogate state with configured
random starts and restart selection. The resulting adversarial images return to
the target branch for adversarial classification and clean/adversarial
embedding consistency.

Recognizer pressure is attached after the target representation is formed.
ArcFace, CosFace, and CurricularFace heads evaluate the same clean or
adversarial samples and contribute weighted losses. The target Partial FC rate
is 0.30, while auxiliary heads use 0.12 to limit classifier-head memory.
Attack generation is chunked to fit the iResNet100 target and Partial FC head
on 24 GB GPUs. The hard-sample term is staged so early epochs learn identity
structure before mined pairs receive stronger weight.

\subsection{Comparison Sequence}
The comparison sequence isolates one source of pressure at a time.
\emph{Baseline-clean} sets the clean recognition reference.
\emph{Recognizer-ensemble-only} tests whether auxiliary recognizer pressure
changes embedding quality without attacks. \emph{Single-attack-only} tests the
risk of narrow adversarial training. \emph{Attack-ensemble-only} measures the
multiattack baseline. \emph{Full-ensemble} then adds recognizer diversity on
top of the same attacker suite, making it the decisive comparison for the main
claim.

\subsection{Threat Model and Evaluation Boundary}
The evaluated attacks are digital image-space attacks, not physical or
presentation attacks. PGD is treated as white-box optimization, BPFA as
transfer-oriented pressure, and DFANet as an FR feature-level attack. During
training, attacks are generated against the current target/surrogate state, so
the defense is not based on hiding the target. Robust classification is
identification-style: it measures whether the closed-set identity prediction
survives perturbation. Pair verification measures a different property:
whether embeddings still separate same-identity from different-identity pairs.

The evaluation is not adaptive. It does not include an adversary that optimizes
through the complete training-time recognizer ensemble, tunes attacks after
observing all defense hyperparameters, or performs query-based score/decision
attacks against the deployed system. The claim is therefore restricted to the
PGD/BPFA/DFANet digital suite and to the ablation question of whether
recognizer-objective pressure adds robustness once attacker diversity is already
present.

\subsection{Experimental Setup and Metrics}
All main experiments use ArcFace/iResNet100 on WebFace4M with 4,015,275
training images and 219,826 validation images over 205,849 identities. Partial
FC uses a 0.30 negative sampling rate to keep large-class training feasible on
8 RTX 4090D GPUs. Main per-GPU batch settings are clean warmup 128, shallow
adversarial 48, full adversarial 32, and attack chunk size 32. Mixed precision,
channels-last tensors, distributed training, and manifest-backed dataloading
are used.

Metrics are grouped by the property they measure. Closed-set validation reports
identity classification. Robust validation reports adversarial accuracy, ASR,
clean-to-robust drop, and per-attack behavior. For attack family $k$,
\begin{equation}
    A^{rob}_{k}=\frac{1}{N}\sum_i
    \mathbf{1}[\hat{y}(x^{adv}_{i,k})=y_i],\quad
    \mathrm{ASR}_k=1-A^{rob}_{k}.
\end{equation}
Verification uses cosine scores
$s_{ij}=g_{\theta}(x_i)^\top g_{\theta}(x_j)$ and reports pair accuracy, EER,
ROC-AUC, TAR at low FAR, and cosine separation. A pair is accepted at
threshold $\tau$ when
\begin{equation}
    d_{\tau}(x_i,x_j)=\mathbf{1}[s_{ij}\geq \tau].
\end{equation}
For genuine pairs $\mathcal{P}^{+}$ and impostor pairs $\mathcal{P}^{-}$,
\begin{equation}
    \mathrm{TAR}(\tau)=\frac{1}{|\mathcal{P}^{+}|}
    \sum_{(i,j)\in\mathcal{P}^{+}} d_{\tau}(x_i,x_j),
\end{equation}
\begin{equation}
    \mathrm{FAR}(\tau)=\frac{1}{|\mathcal{P}^{-}|}
    \sum_{(i,j)\in\mathcal{P}^{-}} d_{\tau}(x_i,x_j).
\end{equation}
EER is the operating point where false accept and false reject rates are equal.
The verification protocol is in-domain because pairs are sampled from the
WebFace4M training manifest and validation identities overlap training
identities. These numbers therefore analyze embedding geometry rather than
external generalization. Pair images are deduplicated, embeddings are extracted
in distributed GPU batches, cosine scores are chunked, and ROC statistics are
computed from sorted scores.
Training cost is reported with wall time, GPU-hours, throughput,
attack-generation share, peak CUDA memory, and active per-GPU batch sizes.
Inference cost is unchanged because auxiliary recognizers and attackers are
discarded after training.

Several implementation choices make the full setting computationally feasible.
Recognizer pressure is stage-gated so clean warmup remains fast. Auxiliary
recognizer heads use a lower Partial FC sampling rate than the target, attack
batches are chunked, verification is batched and distributed, and the
full-adversarial batch size is kept at 32 after larger batch trials increased
memory pressure without improving stable throughput. Training attackers and
evaluation attackers are separated in configuration. Clean-only and
recognizer-only conditions receive no adversarial training pressure, but the
final robustness table evaluates all conditions against the same
PGD/BPFA/DFANet suite.

The evaluation suite remains narrower than the broader deployment threat
model. Stronger AutoAttack/APGD-style checks \cite{autoattack},
TransMix-style transfer attacks \cite{transmix}, certified FR robustness
tests \cite{certrobustfr}, AdvFaceGAN, Adv-Makeup, Greedy-DiM morphing,
diffusion-based attacks, and physical or PAD-aware stress tests are not counted
as evidence here because they require faithful separate implementations or
generated-image protocols
\cite{advfacegan,advmakeup,greedydim,physicalsurvey,fasattack}.

\section{Results}
\subsection{Overview}
The results support a conditional claim rather than a simple ranking.
Recognizer pressure alone preserves clean in-domain verification but does not
make the target robust. Narrow PGD-only pressure improves some adversarial
behavior but damages both classification and verification. The full method
becomes useful only after the attacker side is already diversified:
recognizer-objective pressure improves robustness beyond the attack ensemble,
while moving the embedding away from the clean-verification optimum.

\subsection{Robustness}

\begin{table}[!t]
\centering
\caption{Closed-set clean and adversarial robustness results. Robust average
is over PGD/BPFA/DFANet unless marked otherwise.}
\label{tab:robust}
\resizebox{\columnwidth}{!}{
\begin{tabular}{lcccc}
\toprule
Condition & Clean Acc. & Robust Acc. & Drop & ASR \\
\midrule
Baseline-clean & 99.23 & 0.21 & 99.02 & 99.79 \\
Recognizer-only & 99.23 & 0.23 & 99.01 & 99.77 \\
Single-attack & 79.69 & 34.51 & 45.18 & 55.09 \\
Attack-ensemble & 93.98 & 52.32 & 41.65 & 44.32 \\
Full-ensemble & 93.40 & \textbf{56.55} & \textbf{36.85} & \textbf{39.45} \\
\bottomrule
\end{tabular}}
\vspace{1mm}
\parbox{\columnwidth}{\centering\footnotesize\setlength{\baselineskip}{8pt}
All numbers are percentages. Robust accuracy and ASR are evaluated over PGD/BPFA/DFANet.}
\end{table}

Table \ref{tab:robust} holds the attacker suite fixed for the key comparison.
Adding recognizer-objective pressure on top of PGD, BPFA, and DFANet raises
robust accuracy from 52.32\% to 56.55\%, lowers ASR from 44.32\% to 39.45\%,
and reduces the clean-to-robust drop from 41.65 to 36.85 points. The remaining
difference is therefore attributable to recognizer diversity rather than attack
coverage.

The control rows clarify the boundary. Clean and recognizer-only models remain
almost fully vulnerable under the three-attack suite, with robust accuracy near
0.2000\% and ASR near 99.80\%. Recognizer pressure alone therefore does not
create adversarial robustness. Single-attack training is also incomplete. Its
three-attack average is 34.51\%, with BPFA at 45.71\% but PGD and DFANet both
near 28.90\%, showing weak transfer across attack families.

The per-attack structure is more informative than the average alone. PGD rises
from 45.45\% to 52.55\%, BPFA from 66.09\% to 69.56\%, and DFANet from
45.43\% to 47.55\%. The largest gain is on PGD, suggesting that recognizer
pressure most strongly affects direct margin optimization. DFANet improves
least, consistent with feature-level FR attacks already operating close to the
embedding objective being regularized.

The clean-accuracy movement shows the cost of that gain. Full-ensemble clean
accuracy is slightly lower than attack-ensemble-only (93.40\% versus
93.98\%), so the recognizer heads do not act as simple clean regularizers.
They appear to move the representation toward perturbation stability. The
single-attack row is more severe: PGD-only training reaches 79.69\% clean
accuracy and 34.51\% average robust accuracy, consistent with over-specialized
adversarial pressure degrading the margin classifier while leaving substantial
cross-attack vulnerability.

The epoch histories suggest that this clean-accuracy cost is schedule
sensitive, but not simply unresolved undertraining. In the full condition,
clean validation recovers from 87.57\% at epoch 16 to 92.37\% at epoch 20 and
93.40\% at epoch 24, while robust accuracy rises from 46.15\% to 55.34\% and
then 56.55\%. The late gains therefore taper after the learning-rate drop.
Longer training or a lower-pressure finetuning stage might recover additional
clean accuracy, but the observed trend does not show that the full gap to
attack-ensemble-only would disappear without changing the schedule.

\subsection{Verification}

\begin{table}[!t]
\centering
\caption{In-domain WebFace4M pair-verification metrics. TAR is reported at low
FAR operating points, following standard FR reporting practice.}
\label{tab:verif}
\resizebox{\columnwidth}{!}{
\begin{tabular}{lccccc}
\toprule
Condition & Verif. Acc. & EER & ROC-AUC & TAR@1e-4 & TAR@1e-5 \\
\midrule
Baseline-clean & 99.96 & 0.0470 & \textbf{99.999} & 99.86 & 99.65 \\
Recognizer-only & \textbf{99.96} & \textbf{0.0410} & 99.999 & \textbf{99.92} & \textbf{99.73} \\
Single-attack & 89.74 & 12.82 & 92.01 & 0.37 & 0.11 \\
Attack-ensemble & 98.33 & 1.78 & 99.75 & 85.02 & 73.38 \\
Full-ensemble & 98.21 & 1.93 & 99.73 & 84.81 & 74.16 \\
\bottomrule
\end{tabular}}
\vspace{1mm}
\parbox{\columnwidth}{\centering\footnotesize\setlength{\baselineskip}{8pt}
All numbers are percentages. Lower EER is better; higher values are better for all other columns.}
\end{table}

Table \ref{tab:verif} measures embedding geometry rather than external
generalization. The verification bundle contains 409,196 WebFace4M pairs from
the training manifest, and validation identities overlap training identities.
Near-ceiling clean verification is therefore expected. The informative signal
is the relative degradation across conditions, especially at low-FAR operating
points and in cosine separation.

The clean and recognizer-only rows provide the control. Both retain about
99.96\% verification accuracy, near-perfect ROC-AUC, and EER below 0.0500\%.
The recognizer-only condition slightly improves EER and low-FAR TAR, indicating
that auxiliary margin heads are not inherently harmful to the embedding space.
Degradation appears when margin diversity is coupled to adversarial image
pressure, where the target must satisfy clean identity separation and
perturbed-feature stability at the same time.

The single-attack row shows the strongest geometric failure: 89.74\%
verification accuracy, 12.82\% EER, 92.01\% ROC-AUC, and only 0.37\%
TAR@1e-4. This low-FAR collapse indicates a score-distribution shift, not only
a loss of class accuracy. For FR, closed-set robust accuracy is therefore
insufficient without verification metrics.

Attack-ensemble training restores much of this structure while providing
robustness: 98.33\% verification accuracy, 1.78\% EER, and 99.75\% ROC-AUC.
The full ensemble then improves robust accuracy at a small but consistent
verification cost: 98.21\% accuracy, 1.93\% EER, and TAR@1e-4 of 84.81\%.
The mean negative cosine increases from 0.0440 to 0.0640, indicating weaker
negative-pair separation. The full model is harder to attack, but part of the
clean verification geometry is compressed by the added invariance pressure.

The verification trend gives a similar schedule signal. Full-ensemble
verification partially recovers late in training, moving from 97.70\% accuracy
and 2.43\% EER at epoch 21 to the final 98.21\% accuracy and 1.93\% EER.
Attack-ensemble-only also improves late, reaching 98.33\% accuracy and
1.78\% EER. This makes the verification loss better interpreted as a
schedule-sensitive tradeoff than as a failed convergence point. A longer
low-learning-rate tail, delayed recognizer pressure, or annealed recognizer
weights could plausibly improve pair geometry, but each option changes the
robustness-cost balance that the current ablation measures.

The full method therefore moves the model along a robustness--verification
frontier. Where adversarial probes are plausible, the 4.23-point robust
accuracy gain and 4.87-point ASR reduction may justify the roughly 0.12-point
verification-accuracy loss and 0.15-point EER increase relative to
attack-ensemble-only. Where clean verification dominates the operating
requirement, the clean or recognizer-only rows remain preferable. This is the
main interpretive result: multi-recognizer multi-attacker training should be
viewed as targeted robustness pressure rather than a replacement for clean FR
training.

\subsection{Training Cost}

\begin{table}[!t]
\centering
\caption{Training cost measured from completed 24-epoch histories on 8 GPUs.
Wall time includes training, validation, verification, and model-saving overhead.
Batch lists the dominant active per-GPU batch size.}
\label{tab:cost}
\resizebox{\columnwidth}{!}{
\begin{tabular}{lcccccc}
\toprule
Condition & Batch & Wall h & GPU h & Img/s & Attack Share & Peak GB \\
\midrule
Baseline-clean & 128 & 12.70 & 101.9 & 262.8 & 0.0000 & 20.30 \\
Recognizer-only & 128 & 26.10 & 208.9 & 128.1 & 0.0000 & 20.60 \\
Single-attack & 48/32 & 37.60 & 300.9 & 89.00 & 55.70 & 20.30 \\
Attack-ensemble & 48/32 & 75.50 & 604.2 & 44.30 & 74.90 & 20.30 \\
Full-ensemble & 48/32 & \textbf{110.1} & \textbf{880.7} & 32.70 & 67.70 & \textbf{22.60} \\
\bottomrule
\end{tabular}}
\vspace{1mm}
\parbox{\columnwidth}{\centering\footnotesize\setlength{\baselineskip}{8pt}
Img/s is end-to-end average throughput over the completed training history. Attack share is attack generation as a fraction of measured train compute.}
\end{table}

Table \ref{tab:cost} reports the price of stronger pressure under a
memory-capped batch schedule. The clean baseline takes 12.70 wall-clock hours.
Recognizer-only training reaches 26.10 hours because auxiliary recognizer
heads add backbone and classifier-head passes even without attacks.
Single-attack training takes 37.60 hours, attack-ensemble training takes
75.50 hours, and the full condition reaches 110.1 hours, or 880.7 GPU-hours.
Wall time combines algorithmic cost with batch-size constraint. Clean and
recognizer-only conditions train at batch 128, while adversarial phases are
capped at 48 and 32. The clean baseline therefore uses up to four times the
full-adversarial batch size for a substantial part of training.

The dominant cost is the adversarial inner loop. Attack-ensemble-only spends
74.90\% of measured train compute on attack generation, and the full ensemble
still spends 67.70\% there despite the added recognizer losses. Each attacked
batch performs repeated target/surrogate forward-backward passes before the
ordinary parameter update; the full condition also evaluates recognizer heads
and samples large-class margin weights. Peak memory reaches 22.60 GB on 24 GB
cards, so the method is limited by both repeated computation and reduced batch
size. The measured bottleneck is repeated adversarial optimization under a
large-class FR head.

The full ensemble gains 4.23 robust-accuracy points over attack-ensemble-only
at the cost of 34.60 additional wall-clock hours and 276.50 GPU-hours. The
marginal robustness gain is therefore expensive. Existing acceleration work
suggests practical directions: recycling gradients or adversarial states as in
Free Adversarial Training \cite{freeadv}, reducing repeated full
backpropagation as in YOPO \cite{yopo}, or using cheaper single-step
adversaries with safeguards against catastrophic overfitting \cite{fastadv}.
For this FR setting, the closest direct ablations are lower-frequency
recognizer pressure, cached or replayed adversarial seeds, cheaper early-epoch
attack curricula, shared surrogate forwards across attacks, and selective
full-depth attack refreshes. These target the measured bottleneck without
changing the multi-recognizer multi-attacker objective.

Schedule changes are another path, but they trade runtime for cleaner
optimization dynamics. Extending the final low-learning-rate phase could test
whether verification continues to recover after robustness has mostly
stabilized. Delaying recognizer pressure until the target has stronger
adversarial features could reduce early embedding distortion. Annealing the
recognizer loss late in training could preserve some robustness while allowing
the target margin head to re-separate clean pairs. These variants are
scientifically useful, but the full run already costs 880.7 GPU-hours, so each
additional schedule test should be treated as a deliberate robustness,
verification, and compute tradeoff rather than a free correction.

\section{Discussion and Conclusion}
The work is best read as an ablation-backed framework result. It does not show
state-of-the-art FR, prove that every ensemble member is beneficial, or
establish deployment robustness against all attacks. It shows a narrower but
defensible effect: when a target already receives multiattack pressure, adding
recognizer-objective diversity improves measured robustness under the same
attacker suite. The single-attack result also supports the need for broader
attack-family training rather than narrow PGD-only pressure.

Several limitations define the boundary of this result. The attacker suite is
explicit but not exhaustive. Verification and closed-set robust classification
answer different questions, so both are required for FR robustness analysis.
The verification protocol is in-domain; external unseen-identity benchmarks
such as LFW, CFP-FP, AgeDB, or IJB-style evaluations are needed before
deployment claims. The recognizer ensemble is training-time pressure, not a
deployment ensemble. Training cost also shapes the design itself:
memory-limited batch sizes, repeated attack forwards, and large-class margin
heads make the full condition much slower than the clean baseline. Finally,
the work does not yet establish robustness to adaptive ensemble attacks,
modern semantic or generative attacks, morphing attacks, corruptions, OOD
appearance variation, or physical/PAD-aware attacks.

\subsection{Conclusion}
This study implements a face-recognition robustness framework that trains one
ArcFace/iResNet100 target under explicit recognizer and attacker ensemble
pressure. The result is an ablation-backed tradeoff rather than a leaderboard
claim. Recognizer-objective diversity adds measured adversarial robustness
beyond attacker diversity, but the gain comes with lower clean-verification
performance and a large training-cost increase. The method is therefore best
understood as a practical training framework for FR hardening: useful when
digital adversarial robustness is part of the operating requirement, less
attractive when clean verification is the only priority, and still bounded by
the attacks and in-domain verification protocol evaluated here.

\begin{thebibliography}{00}
\bibitem{deepface} Y. Taigman, M. Yang, M. Ranzato, and L. Wolf,
``DeepFace: Closing the gap to human-level performance in face verification,''
in \emph{Proc. IEEE Conf. Comput. Vis. Pattern Recognit.}, 2014.

\bibitem{facenet} F. Schroff, D. Kalenichenko, and J. Philbin, ``FaceNet: A
unified embedding for face recognition and clustering,'' in \emph{Proc.
IEEE/CVF Conf. Comput. Vis. Pattern Recognit.}, 2015.

\bibitem{vggface} O. M. Parkhi, A. Vedaldi, and A. Zisserman, ``Deep face
recognition,'' in \emph{Proc. Brit. Mach. Vis. Conf.}, 2015.

\bibitem{sphereface} W. Liu, Y. Wen, Z. Yu, M. Li, B. Raj, and L. Song,
``SphereFace: Deep hypersphere embedding for face recognition,'' in
\emph{Proc. IEEE/CVF Conf. Comput. Vis. Pattern Recognit.}, 2017.

\bibitem{arcface} J. Deng, J. Guo, N. Xue, and S. Zafeiriou, ``ArcFace:
Additive angular margin loss for deep face recognition,'' in \emph{Proc.
IEEE/CVF Conf. Comput. Vis. Pattern Recognit.}, 2019, pp. 4690--4699.

\bibitem{cosface} H. Wang, Y. Wang, Z. Zhou, X. Ji, D. Gong, J. Zhou, Z. Li,
and W. Liu, ``CosFace: Large margin cosine loss for deep face recognition,''
in \emph{Proc. IEEE/CVF Conf. Comput. Vis. Pattern Recognit.}, 2018.

\bibitem{curricularface} Y. Huang, Y. Wang, Y. Tai, X. Liu, P. Shen, S. Li,
J. Li, and F. Huang, ``CurricularFace: Adaptive curriculum learning loss for
deep face recognition,'' in \emph{Proc. IEEE/CVF Conf. Comput. Vis. Pattern
Recognit.}, 2020.

\bibitem{magface} Q. Meng, S. Zhao, Z. Huang, and F. Zhou, ``MagFace: A
universal representation for face recognition and quality assessment,'' in
\emph{Proc. IEEE/CVF Conf. Comput. Vis. Pattern Recognit.}, 2021.

\bibitem{adaface} M. Kim, A. K. Jain, and X. Liu, ``AdaFace: Quality adaptive
margin for face recognition,'' in \emph{Proc. IEEE/CVF Conf. Comput. Vis.
Pattern Recognit.}, 2022.

\bibitem{elasticface} F. Boutros, N. Damer, F. Kirchbuchner, and A. Kuijper,
``ElasticFace: Elastic margin loss for deep face recognition,'' in
\emph{Proc. IEEE/CVF Conf. Comput. Vis. Pattern Recognit. Workshops}, 2022.

\bibitem{lfw} G. B. Huang, M. Ramesh, T. Berg, and E. Learned-Miller,
``Labeled Faces in the Wild: A database for studying face recognition in
unconstrained environments,'' Univ. Massachusetts Amherst, Tech. Rep. 07-49,
2007.

\bibitem{megaface} I. Kemelmacher-Shlizerman, S. M. Seitz, D. Miller, and
E. Brossard, ``The MegaFace benchmark: 1 million faces for recognition at
scale,'' in \emph{Proc. IEEE/CVF Conf. Comput. Vis. Pattern Recognit.}, 2016.

\bibitem{ijbc} B. Maze \emph{et al.}, ``IARPA Janus Benchmark-C: Face dataset
and protocol,'' in \emph{Proc. Int. Conf. Biometrics}, 2018.

\bibitem{retinaface} J. Deng, J. Guo, E. Ververas, I. Kotsia, and
S. Zafeiriou, ``RetinaFace: Single-shot multi-level face localisation in the
wild,'' in \emph{Proc. IEEE/CVF Conf. Comput. Vis. Pattern Recognit.}, 2020.

\bibitem{webface260m} Z. Zhu \emph{et al.}, ``WebFace260M: A benchmark
unveiling the power of million-scale deep face recognition,'' in \emph{Proc.
IEEE/CVF Conf. Comput. Vis. Pattern Recognit.}, 2021.

\bibitem{partialfc} X. An \emph{et al.}, ``Partial FC: Training 10 million
identities on a single machine,'' arXiv:2010.05222, 2020.

\bibitem{subcenter} J. Deng, J. Guo, T. Liu, M. Gong, and S. Zafeiriou,
``Sub-center ArcFace: Boosting face recognition by large-scale noisy web
faces,'' in \emph{Proc. Eur. Conf. Comput. Vis.}, 2020.

\bibitem{madry} A. Madry, A. Makelov, L. Schmidt, D. Tsipras, and A. Vladu,
``Towards deep learning models resistant to adversarial attacks,'' in
\emph{Proc. Int. Conf. Learn. Represent.}, 2018.

\bibitem{tramer} F. Tram\`er, A. Kurakin, N. Papernot, I. Goodfellow,
D. Boneh, and P. McDaniel, ``Ensemble adversarial training: Attacks and
defenses,'' in \emph{Proc. Int. Conf. Learn. Represent.}, 2018.

\bibitem{autoattack} F. Croce and M. Hein, ``Reliable evaluation of
adversarial robustness with an ensemble of diverse parameter-free attacks,''
in \emph{Proc. Int. Conf. Mach. Learn.}, 2020.

\bibitem{dongattack} Y. Dong, H. Su, B. Wu, Z. Li, W. Liu, T. Zhang, and
J. Zhu, ``Efficient decision-based black-box adversarial attacks on face
recognition,'' in \emph{Proc. IEEE/CVF Conf. Comput. Vis. Pattern Recognit.},
2019, pp. 7714--7722.

\bibitem{frsurvey} F. Vakhshiteh, A. Nickabadi, and R. Ramachandra,
``Adversarial attacks against face recognition: A comprehensive study,''
\emph{IEEE Access}, vol. 9, pp. 92735--92756, 2021.

\bibitem{facesec} Z. Jiang, T. Chen, T. Chen, and Z. Wang, ``FACESEC: A
fine-grained robustness evaluation framework for face recognition systems,''
in \emph{Proc. IEEE/CVF Conf. Comput. Vis. Pattern Recognit. Workshops}, 2021.

\bibitem{robustbench} F. Croce, M. Andriushchenko, V. Sehwag, E. Debenedetti,
N. Flammarion, M. Chiang, P. Mittal, and M. Hein, ``RobustBench: A
standardized adversarial robustness benchmark,'' in \emph{Adv. Neural Inf.
Process. Syst.}, 2021.

\bibitem{multirobustbench} S. Dai, S. Mahloujifar, C. Xiang, V. Sehwag,
P.-Y. Chen, and P. Mittal, ``MultiRobustBench: Benchmarking robustness against
multiple attacks,'' in \emph{Proc. Int. Conf. Mach. Learn.}, 2023.

\bibitem{oodface} C. Kang, Y. Chen, S. Ruan, S. Zhao, R. Zhang, J. Wang,
S. Fu, and X. Wei, ``OODFace: Benchmarking robustness of face recognition under
common corruptions and appearance variations,'' arXiv:2412.02479, 2024.

\bibitem{freeadv} A. Shafahi, M. Najibi, M. A. Ghiasi, Z. Xu, J. Dickerson,
C. Studer, L. S. Davis, G. Taylor, and T. Goldstein, ``Adversarial training
for free!'' in \emph{Adv. Neural Inf. Process. Syst.}, 2019.

\bibitem{yopo} D. Zhang, T. Zhang, Y. Lu, Z. Zhu, and B. Dong, ``You only
propagate once: Accelerating adversarial training via maximal principle,'' in
\emph{Adv. Neural Inf. Process. Syst.}, 2019.

\bibitem{fastadv} E. Wong, L. Rice, and J. Z. Kolter, ``Fast is better than
free: Revisiting adversarial training,'' in \emph{Proc. Int. Conf. Learn.
Represent.}, 2020.

\bibitem{transmix} Y. M. Khedr, X. Liu, and K. He, ``TransMix: Crafting
highly transferable adversarial examples to evade face recognition models,''
\emph{Image Vis. Comput.}, vol. 146, Art. no. 105022, 2024.

\bibitem{certrobustfr} S. Paik, D. Kim, C. Hwang, S. Kim, and J. H. Seo,
``Towards certifiably robust face recognition,'' in \emph{Proc. Eur. Conf.
Comput. Vis.}, 2024.

\bibitem{physicalsurvey} M. Wang, J. Zhou, T. Li, G. Meng, and K. Chen, ``A
survey on physical adversarial attacks against face recognition systems,''
\emph{Neurocomputing}, 2025.

\bibitem{advmakeup} B. Yin, W. Wang, T. Yao, J. Guo, Z. Kong, S. Ding,
J. Li, and C. Liu, ``Adv-Makeup: A new imperceptible and transferable attack
on face recognition,'' in \emph{Proc. Int. Joint Conf. Artif. Intell.}, 2021,
pp. 1252--1258.

\bibitem{greedydim} Z. W. Blasingame and C. Liu, ``Greedy-DiM: Greedy
algorithms for unreasonably effective face morphs,'' in \emph{Proc. IEEE Int.
Joint Conf. Biometrics}, 2024.

\bibitem{advfacegan} H. Huang, Y. Yang, and Y. Wang, ``AdvFaceGAN: A face
dual-identity impersonation attack method based on generative adversarial
networks,'' \emph{PeerJ Computer Science}, vol. 11, e2904, 2025.

\bibitem{fasattack} X. Dong, J. Zhu, H. Su, D. Wang, W. Zhang, and Y. Yu,
``Adversarial attacks on both face recognition and face anti-spoofing models,''
in \emph{Proc. Int. Joint Conf. Artif. Intell.}, 2025.
\end{thebibliography}

\end{document}
